\songtitle{The Long Way to Paris}

\tag*{2026}\tag{hip-hop}\label{sec:long-way-to-paris}
\begin{notes}
Lyrics by Claude, written in an attempt to mimic Carlos Thompson's voice and songwriting style following specific prompts about the overall subject: an alternate-history legend of a Haitian general who takes Paris and weds the Republic's Crown Princess. Retrieved from Gabriel's Suno page (V6-MINI, 17 October 2026).
\end{notes}
\begin{style}
Historical hip-hop ballad, 84 BPM half-time march, booming drums with orchestral weight, deep-baritone spoken-word verses rising into powerful full-voice belts and gang-vocal shouts on the hook; orchestral string ostinato, taiko-style war drums, djembé and kora motifs, brass fanfare turning tender in the last chorus.
\end{style}
\begin{lyrics}
\annotation{Spoken intro}\\
This is the story of a general.\\
He was not born to a crown.\\
He was born where the cane was cut,\\
and he learned arithmetic from the ledger that listed his name.

\songpart{Verse 1}
Once a fleet crossed the ocean to take the island back.\\
The island answered with fever, fire and the mountain road.\\
Thirty years later the ships went the other way:\\
forty vessels from Cap-Haïtien, a general at the bow.

Brest in a week. Rouen in a month.\\
Paris listened for the drums and asked what he wanted.\\
He answered in the language of treaties:
"Everything. In writing."

\songpart{Chorus}
General of the island, conqueror of France,\\
he took the capital by autumn\\
and the republic by treaty.\\
As told by the legend.

\songpart{Verse 2}
The Republic had abolished the crown in a very serious year.\\
It kept a Crown Princess anyway, for the tourists and the paperwork.\\
Nobody could explain it. The constitution had a footnote.\\
The general read the footnote twice.\\
She had written it herself.

\songpart{Verse 3}
The Crown Princess said: you took my city, my ministers, my bread.\\
The general said: I took nothing I can't return.\\
Case one: a marriage as hostage. Reject.\\
Case two: a marriage as treaty. Accept.\\
Case three: a marriage as marriage. Unclear. Proceed.

\songpart{Chorus}
General of the island, conqueror of France,\\
he took the capital by autumn\\
and the republic by wedding.\\
As told by the legend.

\songpart{Bridge, spoken}
Article one: the ledgers burn, and the names are read aloud first.\\
Article two: the Republic pays for what the Republic took,\\
in coin, and in schools.\\
Article three: the Crown Princess keeps the title,\\
since nobody remembers how to revoke it.\\
Article four: the wedding is Saturday.

\songpart{Verse 4}
The bells of Paris rang the bells that once rang for the fleet.\\
Nobody mentioned it. Everybody heard.\\
You saved the city, said the Crown Princess.\\
You saved the general, said the city.\\
Nou la.

\songpart{Final Chorus}
General of the island, conqueror of France,\\
he took the capital by autumn\\
and the republic, mostly, by being decent.\\
As told by the legend.

\songpart{Outro, spoken, instruments out}
They grew old in the palace gardens\\
and planted coffee in a greenhouse that never liked it.\\
The coffee failed.\\
The marriage held.\\
The legend is a little exaggerated.\\
The marriage held.
\end{lyrics}

